\chapitrepage{Conception du système}

\begin{objectifbox}
Ce chapitre répond à la question « comment construire la solution ? ». Nous
présentons les scénarios dynamiques, l'architecture en couches, les modèles de
prévision, la conception du MRP, le modèle logique des données et le déploiement.
\end{objectifbox}

\section{Introduction}

La conception de \NomProjet repose sur une idée centrale : séparer la donnée
source, les calculs métier et la restitution. Cette séparation permet de rejouer
un calcul, de tester chaque règle indépendamment et de remplacer un composant sans
réécrire toute l'application. Les fichiers CSV et JSON forment un contrat simple
entre les modules ; le dashboard n'entraîne pas lui-même les modèles, il lit les
résultats produits par les moteurs.

\section{Modélisation dynamique}

\subsection{Séquence de mise à jour quotidienne}

La figure~\ref{fig:sequence-maj} détaille le scénario automatisé. L'extraction SQL
est volontairement placée avant les deux moteurs. Les fichiers journaliers sont
générés avant le pipeline mensuel, car ce dernier peut compléter son panel avec les
mois complets issus de la source journalière.

\begin{figure}[H]
\centering
\resizebox{\textwidth}{!}{%
\begin{tikzpicture}[>=Latex,font=\scriptsize]
  \foreach \x/\name/\label in {
    0/task/Planificateur,
    3/sql/SQL Elyx,
    6/orch/Orchestrateur,
    9/day/Moteur journalier,
    12/month/Pipeline mensuel,
    15/files/Exports}
  {
    \node[draw,rounded corners,fill=paulCream,minimum width=2.35cm,minimum height=0.75cm] (\name) at (\x,0) {\label};
    \draw[dashed,gray] (\x,-0.4) -- (\x,-8.1);
  }
  \draw[->] (task.south) ++(0,-0.65) -- node[above]{déclencher à 05 h 30} ($(orch.south)+(0,-0.65)$);
  \draw[->] ($(orch.south)+(0,-1.45)$) -- node[above]{lire jours clôturés} ($(sql.south)+(0,-1.45)$);
  \draw[->] ($(sql.south)+(0,-2.15)$) -- node[above]{lignes agrégées} ($(orch.south)+(0,-2.15)$);
  \draw[->] ($(orch.south)+(0,-2.85)$) -- node[above]{fusion contrôlée CSV} ($(files.south)+(0,-2.85)$);
  \draw[->] ($(orch.south)+(0,-3.65)$) -- node[above]{générer 62 jours} ($(day.south)+(0,-3.65)$);
  \draw[->] ($(day.south)+(0,-4.35)$) -- node[above]{prévisions + suivi} ($(files.south)+(0,-4.35)$);
  \draw[->] ($(orch.south)+(0,-5.15)$) -- node[above]{lancer le pipeline} ($(month.south)+(0,-5.15)$);
  \draw[->] ($(month.south)+(0,-5.85)$) -- node[above]{plan + MRP + validation} ($(files.south)+(0,-5.85)$);
  \draw[->] ($(files.south)+(0,-6.65)$) -- node[above]{statut et journaux} ($(orch.south)+(0,-6.65)$);
  \draw[->] ($(orch.south)+(0,-7.35)$) -- node[above]{code retour} ($(task.south)+(0,-7.35)$);
\end{tikzpicture}}
\caption{Diagramme de séquence de la mise à jour quotidienne}
\label{fig:sequence-maj}
\end{figure}

\subsection{Activité de génération d'une recommandation}

Le calcul d'une recommandation suit une série de contrôles. Une vente future ne
doit jamais être mélangée aux données d'entraînement et un produit sans historique
suffisant doit être signalé au lieu de recevoir une fausse précision.

\begin{figure}[H]
\centering
\resizebox{0.98\textwidth}{!}{%
\begin{tikzpicture}[
  node distance=8mm and 14mm,
  flowstep/.style={draw,rounded corners,fill=white,minimum width=6.2cm,minimum height=0.72cm,align=center,font=\small},
  decision/.style={draw,diamond,aspect=2.6,fill=paulCream,align=center,font=\small},
  arrow/.style={-{Latex[length=2mm]},thick,paulBrown}]
  \node[flowstep] (load) {Charger et nettoyer la série du produit};
  \node[decision,below=of load] (hist) {Historique suffisant ?};
  \node[flowstep,below left=of hist] (fallback) {Appliquer un repli prudent\\et marquer « historique court »};
  \node[flowstep,below right=of hist] (model) {Calculer le modèle sélectionné\\et les composantes calendaires};
  \node[flowstep,below=18mm of hist] (boost) {Combiner fêtes, matchs, événements et commandes};
  \node[flowstep,below=of boost] (unc) {Estimer l'incertitude et la marge de sécurité};
  \node[flowstep,below=of unc] (override) {Appliquer une correction manuelle éventuelle};
  \node[flowstep,below=of override] (out) {Écrire prévision, recommandation et fiabilité};
  \draw[arrow] (load) -- (hist);
  \draw[arrow] (hist) -- node[above left]{non} (fallback);
  \draw[arrow] (hist) -- node[above right]{oui} (model);
  \draw[arrow] (fallback) |- (boost);
  \draw[arrow] (model) |- (boost);
  \draw[arrow] (boost) -- (unc);
  \draw[arrow] (unc) -- (override);
  \draw[arrow] (override) -- (out);
\end{tikzpicture}
}
\caption{Diagramme d'activité du calcul d'une recommandation}
\label{fig:activite-recommandation}
\end{figure}

\section{Architecture logicielle}

\subsection{Architecture globale}

L'architecture de la figure~\ref{fig:architecture-globale} distingue six zones :
sources, ingestion, moteurs, règles métier, sorties et interface. Les flèches
pleines représentent le flux nominal ; la configuration JSON agit transversalement
sur les moteurs et les règles.

\begin{figure}[H]
\centering
\resizebox{\textwidth}{!}{%
\begin{tikzpicture}[
  box/.style={draw,rounded corners,minimum width=3.2cm,minimum height=1.0cm,align=center,font=\small},
  source/.style={box,fill=softGray},
  engine/.style={box,fill=paulCream},
  output/.style={box,fill=white},
  flow/.style={-{Latex[length=2.2mm]},thick,paulBrown}]

  \node[source] (dos) at (0,3.8) {Exports DOS / Excel\\historiques};
  \node[source] (sql) at (0,2.1) {SQL Server Elyx\\jours clôturés};
  \node[source] (refs) at (0,0.4) {Recettes, prix, fêtes,\\événements, commandes};

  \node[engine] (ingest) at (4.2,3.0) {Ingestion et nettoyage\\\texttt{data\_loader.py}};
  \node[engine] (config) at (4.2,0.8) {Configuration métier\\\texttt{config.py} + JSON};

  \node[engine] (daily) at (8.5,4.1) {Prévision journalière\\62 jours};
  \node[engine] (monthly) at (8.5,2.5) {Prévision mensuelle\\12 mois};
  \node[engine] (rules) at (8.5,0.9) {Fêtes, matchs, commandes,\\incertitude};
  \node[engine] (mrp) at (8.5,-0.7) {BOM / PSF / coûts\\et couverture};

  \node[output] (csv) at (12.8,3.6) {CSV et Excel\\opérationnels};
  \node[output] (reports) at (12.8,1.8) {Rapports, graphiques\\et journaux};
  \node[output] (dash) at (12.8,0) {Dashboard Dash\\et assistant local};

  \draw[flow] (dos) -- (ingest);
  \draw[flow] (sql) -- (ingest);
  \draw[flow] (refs) -- (config);
  \draw[flow] (ingest) -- (daily);
  \draw[flow] (ingest) -- (monthly);
  \draw[flow] (config) -- (rules);
  \draw[flow] (config) -- (mrp);
  \draw[flow] (daily) -- (csv);
  \draw[flow] (monthly) -- (csv);
  \draw[flow] (rules) -- (daily);
  \draw[flow] (rules) -- (monthly);
  \draw[flow] (monthly) -- (mrp);
  \draw[flow] (mrp) -- (reports);
  \draw[flow] (csv) -- (dash);
  \draw[flow] (reports) -- (dash);

  \node[draw,dashed,rounded corners,fit=(ingest)(mrp),inner sep=5mm] (calcbox) {};
  \node[font=\bfseries,anchor=south west,xshift=2mm,yshift=2mm]
        at (calcbox.north west) {Couche de calcul Python};
\end{tikzpicture}}
\caption{Architecture fonctionnelle globale de PrevisionPaul}
\label{fig:architecture-globale}
\end{figure}

\subsection{Responsabilités des modules}

Le code principal est réparti dans le paquet \texttt{paul\_forecast}. Le tableau
\ref{tab:modules} synthétise les responsabilités. Cette organisation joue le rôle
d'un diagramme de composants statique : les dépendances vont de la configuration
vers les services métier, puis vers le pipeline et l'interface.

\small
\begin{longtable}{p{4.5cm}p{9.9cm}}
\caption{Principaux composants logiciels}\label{tab:modules}\\
\toprule
\textbf{Composant} & \textbf{Responsabilité} \\
\midrule
\endfirsthead
\toprule
\textbf{Composant} & \textbf{Responsabilité} \\
\midrule
\endhead
\texttt{data\_loader.py} & Charge, valide, nettoie et agrège les ventes mensuelles et journalières. \\
\texttt{forecasting.py} & Implémente les prévisions de base et les replis lorsque certaines bibliothèques sont absentes. \\
\texttt{modele\_global.py} & Entraîne un modèle de gradient boosting partagé entre les produits à partir de variables retardées et calendaires. \\
\texttt{forecast\_journalier.py} & Calcule le niveau récent, l'ancre annuelle, le profil hebdomadaire, les ruptures et l'horizon de 62 jours. \\
\texttt{saisonnalite\_fetes.py} / \texttt{calibration\_fetes.py} & Mesure et applique les profils des fêtes marocaines. \\
\texttt{matchs.py} / \texttt{evenements.py} & Apprend ou applique les impacts des matchs et événements ponctuels. \\
\texttt{commandes.py} & Détecte les pics assimilables à des commandes, nettoie l'apprentissage et ajoute les commandes futures. \\
\texttt{incertitude.py} & Convertit les erreurs mesurées en intervalles, quantités recommandées et étiquettes de fiabilité. \\
\texttt{bom.py} & Associe les produits aux recettes, éclate les PSF et normalise les matières. \\
\texttt{couts.py} / \texttt{marges.py} & Valorise les besoins et calcule les indicateurs de marge matière. \\
\texttt{couverture.py} & Mesure la part des volumes couverte par des recettes exactes ou estimées. \\
\texttt{backtest.py} / \texttt{suivi.py} & Évalue les modèles sans fuite temporelle aux niveaux mensuel et journalier. \\
\texttt{pipeline.py} & Orchestre la prévision mensuelle, la réconciliation, le MRP, la validation et les exports. \\
\texttt{assistant.py} & Fournit des outils de question-réponse déterministes à partir des fichiers locaux. \\
\texttt{reporting.py} & Génère tableaux, classeurs, bons de commande et graphiques statiques. \\
\texttt{dashboard.py} & Construit l'interface Dash, ses vues, ses callbacks et ses exports interactifs. \\
\bottomrule
\end{longtable}
\normalsize

\section{Conception des moteurs de prévision}

\subsection{Sélection mensuelle par produit}

Un seul modèle global appliqué à toutes les références serait inadapté. Nous
comparons sept modèles de base et trois candidats complémentaires, soit dix
candidats : ETS/Holt--Winters, Theta, ARIMA, naïf saisonnier, moyenne mobile,
moyenne historique, Croston, ensemble moyen, ensemble médian et gradient boosting
global. Ces familles sont décrites dans la littérature de référence sur la
prévision~\cite{hyndman2021} et dans les documentations de
statsmodels~\cite{statsmodelsHW} et scikit-learn~\cite{sklearnGBM}. Croston cible
notamment les demandes intermittentes ; le modèle naïf
saisonnier sert de référence ; les ensembles réduisent la sensibilité à un choix
unique ; le gradient boosting mutualise l'information entre produits.

\begin{table}[H]
\centering
\caption{Candidats de prévision mensuelle}
\label{tab:modeles}
\small
\begin{tabularx}{\textwidth}{p{3.1cm}Yp{4.2cm}}
\toprule
\textbf{Candidat} & \textbf{Principe} & \textbf{Profil adapté} \\
\midrule
ETS / Holt--Winters & Niveau, tendance et saisonnalité exponentiels. & Séries régulières avec saisonnalité annuelle. \\
Theta & Décomposition de la série en lignes Theta. & Séries lisses ou modérément tendancielles. \\
ARIMA & Dépendances autorégressives et moyennes mobiles. & Séries avec autocorrélation exploitable. \\
Naïf saisonnier & Reprend la valeur du même mois de l'année précédente. & Saisonnalité stable et forte. \\
Moyenne mobile & Moyenne des derniers mois. & Niveau récent stable. \\
Moyenne & Moyenne historique. & Référence simple ou historique court. \\
Croston & Sépare taille et intervalle des demandes non nulles. & Produits intermittents. \\
Ensemble moyen & Moyenne de plusieurs modèles statistiques. & Réduction du risque d'un modèle isolé. \\
Ensemble médian & Médiane robuste de quatre candidats. & Présence de prévisions extrêmes. \\
GBM global & Apprentissage partagé sur retards, calendrier, fêtes et famille. & Produits bénéficiant de motifs communs. \\
\bottomrule
\end{tabularx}
\end{table}

Le benchmark utilise douze fenêtres glissantes à un mois d'horizon. Pour chaque
produit $p$ et fenêtre $t$, l'erreur absolue est :

\begin{equation}
e_{p,t,m}=\left|y_{p,t}-\widehat{y}_{p,t,m}\right|,
\end{equation}

où $m$ désigne le modèle candidat. Le modèle retenu minimise la moyenne des
erreurs disponibles. Le fichier \texttt{modele\_par\_produit.json} conserve ce
choix pour 848 produits disposant d'un historique suffisant.

\subsection{Réconciliation hiérarchique}

La somme de plus de mille prévisions ajustées séparément peut diverger du niveau
global. Nous calculons donc une prévision agrégée Holt--Winters, puis un facteur de
réconciliation mensuel :

\begin{equation}
r_t = \mathrm{clip}\left(
\frac{\widehat{Y}^{\,\mathrm{global}}_t}
{\sum_p \widehat{y}_{p,t}},\;0{,}7,\;1{,}4\right).
\end{equation}

La prévision réconciliée devient
$\widehat{y}^{\,\mathrm{rec}}_{p,t}=r_t\widehat{y}_{p,t}$. Les bornes évitent qu'un
agrégat aberrant ne transforme excessivement toutes les références.

\subsection{Modèle journalier}

Le moteur journalier cherche un compromis entre réactivité et saisonnalité. Pour
un jour futur $d$, il calcule :

\begin{equation}
\widehat{y}_{p,d}=\left[\beta_p L^{\,\text{récent}}_p+
(1-\beta_p)L^{\,\text{annuel}}_{p,d}\right]
\times w_{p,j(d)} \times b_{p,d},
\end{equation}

où $w_{p,j(d)}$ est le poids du jour de semaine et $b_{p,d}$ le multiplicateur
métier combiné. En régime normal, $\beta=0{,}4$. Lorsqu'une hausse ou une baisse
durable est détectée par comparaison de fenêtres récentes avec l'ancre annuelle,
le poids récent est porté jusqu'à $0{,}85$ ou $0{,}90$. Le ratio de croissance
annuel est borné entre 0,6 et 1,8 afin de contenir les valeurs extrêmes.

La fiabilité journalière est mesurée en rejouant les 28 derniers jours. Les
paramètres sont réestimés chaque semaine et l'erreur porte sur les totaux
hebdomadaires. Le produit est classé « Fiable », « Moyen », « Incertain », « Peu
vendu » ou « Historique court ». Ces libellés évitent de présenter un pourcentage
fragile pour une référence presque jamais vendue.

\subsection{Combinaison des événements et commandes}

Les multiplicateurs calendaires ne sont pas simplement multipliés les uns par les
autres. Le système conserve le plus grand facteur supérieur à un et le plus petit
facteur inférieur à un, puis les combine. Cette règle empêche qu'une fête, un
match et un événement saisi le même jour produisent une hausse irréaliste.

Les commandes clients suivent une règle différente. Une commande future connue
est une quantité certaine : elle est ajoutée à la prévision de la date et au MRP.
En revanche, un pic historique détecté comme atypique est remplacé par une médiane
locale pendant l'apprentissage. Cette séparation évite que le modèle interprète
une commande exceptionnelle comme une nouvelle demande habituelle.

\section{Conception du calcul MRP}

\subsection{Explosion des nomenclatures}

Pour un produit $p$, une matière $i$ et un mois $t$, le besoin brut est :

\begin{equation}
B_{i,t}=\sum_p Q^{\,\text{recommandée}}_{p,t}\times q_{p,i},
\end{equation}

où $q_{p,i}$ représente la quantité de matière $i$ nécessaire pour une unité du
produit. Lorsque la recette contient un produit semi-fini, celui-ci est remplacé
récursivement par sa propre recette. Un ensemble d'alias unifie par exemple les
différentes écritures d'une même farine. Les grammes et millilitres sont ensuite
convertis en kilogrammes et litres pour le bon de commande.

L'algorithme applique l'ordre de priorité suivant : recette exacte validée,
détection par le nom, recette générique de famille, puis absence explicite de
recette. La provenance accompagne le résultat. Ainsi, une estimation reste
utilisable pour le pilotage, mais elle n'est pas confondue avec une fiche validée
par le chef.

\subsection{Valorisation}

Les prix sont recherchés par mots-clés du plus spécifique au plus général. Le coût
d'une ligne vaut la quantité normalisée multipliée par le prix unitaire. Le
food-cost présenté dans le dashboard est un plancher, car il ne comprend ni
main-d'œuvre, ni énergie, ni emballages absents du référentiel. Cette restriction
est affichée à l'utilisateur au lieu d'être masquée.

\section{Modèle logique des données}

\subsection{Relations principales}

La solution n'impose pas une base applicative transactionnelle. Les tables
logiques sont matérialisées par des CSV et JSON. La figure~\ref{fig:modele-donnees}
montre les relations nécessaires au calcul. Les transformations tabulaires et
temporelles s'appuient principalement sur pandas~\cite{pandasDocs}.

\begin{figure}[H]
\centering
\resizebox{0.96\textwidth}{!}{%
\begin{tikzpicture}[
  entity/.style={draw,rounded corners,fill=white,minimum width=4.2cm,align=left,font=\scriptsize,inner sep=3mm},
  rel/.style={-{Latex[length=2mm]},thick,paulBrown}]
  \node[entity] (vente) at (0,4.2) {\textbf{VENTE JOURNALIÈRE}\\Date, Code, Produit, Famille\\Quantité, CA TTC};
  \node[entity] (produit) at (6,4.2) {\textbf{PRODUIT}\\Nom, Famille, Modèle retenu\\Ajustement, provenance recette};
  \node[entity] (recette) at (12,4.2) {\textbf{RECETTE}\\Produit, Ingrédient\\Quantité unitaire, source};
  \node[entity] (cal) at (0,0.8) {\textbf{CALENDRIER MÉTIER}\\Fête, événement, match\\Période, portée, impact};
  \node[entity] (prev) at (6,0.8) {\textbf{PRÉVISION}\\Date/Mois, Produit, Quantité\\Recommandation, fiabilité};
  \node[entity] (mat) at (12,0.8) {\textbf{BESOIN MATIÈRE}\\Date/Mois, Ingrédient\\Quantité, unité, coût};
  \node[entity] (cmd) at (0,-2.5) {\textbf{COMMANDE CLIENT}\\Date, Produit, Quantité\\Client, identifiant};
  \node[entity] (prix) at (12,-2.5) {\textbf{PRIX MATIÈRE}\\Mot-clé, prix, unité\\Devise, date de mise à jour};

  \draw[rel] (vente) -- node[above]{historise} (produit);
  \draw[rel] (produit) -- node[above]{possède} (recette);
  \draw[rel] (vente) -- (prev);
  \draw[rel] (cal) -- (prev);
  \draw[rel] (cmd) -- (prev);
  \draw[rel] (produit) -- (prev);
  \draw[rel] (prev) -- (mat);
  \draw[rel] (recette) -- (mat);
  \draw[rel] (prix) -- (mat);
\end{tikzpicture}}
\caption{Modèle logique des données de PrevisionPaul}
\label{fig:modele-donnees}
\end{figure}

\subsection{Dictionnaire synthétique}

\begin{table}[H]
\centering
\caption{Dictionnaire des champs principaux}
\label{tab:dictionnaire}
\small
\begin{tabularx}{\textwidth}{p{3.2cm}p{2.6cm}Yp{3.2cm}}
\toprule
\textbf{Champ} & \textbf{Type} & \textbf{Définition} & \textbf{Contrainte} \\
\midrule
\texttt{Date} & date ISO & Jour de vente ou de prévision. & Jour clôturé pour les ventes SQL. \\
\texttt{Produit} & texte & Libellé canonique de la référence. & Non vide ; rapproché des alias. \\
\texttt{Famille} & catégorie & Département commercial. & Une des familles connues. \\
\texttt{Quantite} & décimal & Quantité réellement vendue. & Valeur négative ramenée ou traitée comme retour. \\
\texttt{Qty\_Prev} & décimal & Prévision centrale avant sécurité. & Supérieure ou égale à zéro. \\
\texttt{Qty\_Recommandee} & décimal & Prévision augmentée selon l'incertitude. & Supérieure ou égale à \texttt{Qty\_Prev}. \\
\texttt{Fiabilite} & catégorie & Niveau issu du backtest. & Fiable, Moyen, Incertain, Peu vendu ou Hist. court. \\
\texttt{Ingredient} & texte & Matière première canonique. & Alias résolus avant agrégation. \\
\texttt{Quantite\_Requise} & décimal & Besoin total calculé. & Unité normalisée pour l'achat. \\
\texttt{Modele} & catégorie & Candidat retenu pour le produit. & Présent dans la liste des candidats. \\
\bottomrule
\end{tabularx}
\end{table}

\section{Architecture matérielle et déploiement}

La solution s'exécute sur un serveur ou un poste Windows relié au réseau interne.
Le serveur SQL Elyx reste la source transactionnelle ; \NomProjet{} effectue une
lecture sur les jours clôturés et conserve son propre historique analytique en
fichiers. Le dashboard écoute par défaut sur le port local 8050. Le navigateur
peut être ouvert sur la même machine, ce qui évite toute exposition Internet.

\begin{figure}[H]
\centering
\resizebox{0.98\textwidth}{!}{%
\begin{tikzpicture}[
  node distance=16mm,
  nodebox/.style={draw,rounded corners,fill=paulCream,minimum width=4.8cm,minimum height=1.2cm,align=center},
  flow/.style={-{Latex[length=2mm]},thick,paulBrown}]
  \node[nodebox] (caisse) {Caisses / Elyx\\SQL Server Express 2014};
  \node[nodebox,right=of caisse] (server) {Serveur Windows\\Python + PrevisionPaul};
  \node[nodebox,right=of server] (browser) {Navigateur local\\Dashboard Dash :8050};
  \node[nodebox,below=of server] (disk) {Stockage local\\CSV, JSON, XLSX, PNG, logs};
  \node[nodebox,above=of server] (task) {Planificateur de tâches\\exécution quotidienne 05 h 30};
  \draw[flow] (caisse) -- node[above]{ODBC / lecture} (server);
  \draw[flow] (task) -- (server);
  \draw[flow] (server) -- (disk);
  \draw[flow] (disk) -- (server);
  \draw[flow] (server) -- node[above]{HTTP local} (browser);
\end{tikzpicture}
}
\caption{Diagramme de déploiement matériel}
\label{fig:deploiement}
\end{figure}

\section{Conclusion}

La conception retenue combine des scénarios automatisés, des moteurs spécialisés
et des contrats de données simples. La validation temporelle, la réconciliation,
la séparation des commandes exceptionnelles et l'explosion récursive des recettes
répondent directement aux contraintes identifiées. Le chapitre suivant montre
comment cette conception a été implémentée, testée et présentée aux utilisateurs.
